\chapter{연산과 배선에서 timing을 예측하는 직관}

\begin{summarybox}[Timing을 읽는 핵심]
연산자마다 고정된 지연 시간이 있는 것이 아니다. 같은 adder도 bit 폭, carry 앞뒤
MUX, source와 destination의 위치, fanout과 congestion에 따라 달라진다. 먼저
register-to-register 경로를 그리고 logic delay와 route delay를 나눈 뒤, clock
budget 안에 들어가는지 판단한다.
\end{summarybox}

\section{주파수를 한 cycle의 시간으로 바꾼다}

clock period는 $T=1/f$다.

\begin{longtable}{C{32mm}C{42mm}L{82mm}}
\toprule
Frequency & Nominal period & 설계자가 기억할 점 \\
\midrule
\endhead
100~MHz & 10.000~ns & 실제 data-path budget은 uncertainty와 setup 때문에 더 작다. \\
130~MHz & 7.692~ns & 8~ns에 가까운 조합 경로는 들어가지 않는다. \\
150~MHz & 6.667~ns & carry 두 번, 큰 MUX와 route가 겹치면 빠듯하다. \\
200~MHz & 5.000~ns & 넓은 arithmetic과 memory control은 적극적인 stage 분할이 필요하다. \\
300~MHz & 3.333~ns & logic뿐 아니라 BRAM/DSP/clock primitive의 물리 한계도 확인한다. \\
\bottomrule
\end{longtable}

setup 조건은 개념적으로 다음과 같다.

\[
T_{clk} \ge
T_{cq}+T_{logic}+T_{route}+T_{setup}+T_{uncertainty}-T_{skew}.
\]

10 ns clock이라고 해서 RTL의 조합 회로에 10 ns를 모두 쓸 수 있는 것이 아니다.
Clock-to-Q, setup, skew, uncertainty가 빠지고, 실제 placement 뒤에는 logic cell
사이 route delay가 더해진다.

\section{Slack과 data path delay}

required time에서 arrival time을 뺀 값이 setup slack이다. Worst Negative Slack이
음수이면 가장 나쁜 setup path가 제약을 넘었다.

\begin{itemize}
  \item \wns $=+0.089$~ns: 최악 경로도 약 0.089 ns의 여유를 가진다.
  \item \wns $=-0.338$~ns: 최악 경로가 요구 시각보다 약 0.338 ns 늦다.
  \item Total Negative Slack은 실패 endpoint 전체의 정도를 함께 보여 준다.
\end{itemize}

작은 양의 WNS는 성공이지만 tool/version/seed/상위 시스템이 바뀌어도 넉넉하다는
뜻은 아니다. timing-closed라는 사실과 margin이 충분하다는 판단을 구분한다.

\section{한 개 LUT가 몇 ns라고 외울 수 없는 이유}

LUT 내부 지연은 선택된 input pin, function과 output 경로에 따라 달라지고, 다음
cell까지의 route는 placement에 따라 달라진다. carry와 MUXF는 일반 routing과
다른 dedicated path를 쓴다. 따라서 ``LUT 하나는 0.2 ns''처럼 고정 상수를
적용하면 실제 critical path를 잘못 예측하기 쉽다.

대신 다음 순서로 빠르게 분류한다.

\begin{enumerate}
  \item LUT level이 한두 단계인지 여러 단계의 decode/MUX tree인지 본다.
  \item CARRY4가 직렬로 몇 개 연결되는지 센다.
  \item DSP/BRAM input 또는 output register가 활성화되어 있는지 본다.
  \item source net의 fanout과 destination까지 physical distance를 본다.
  \item post-route report의 logic/route 비율로 다음 수정 방향을 정한다.
\end{enumerate}

\section{실측 경로로 얻은 ns 범위}

다음 값은 같은 Unified NTT 개발 과정에서 관찰한 post-route data path다. 특정
operator의 보편적 delay가 아니라, 어떤 구조 조합이 어느 정도 budget을 요구했는지
보여 주는 경험 범위다.

\begin{longtable}{L{45mm}C{22mm}C{22mm}C{22mm}L{32mm}}
\toprule
경로 성격 & Data path & Logic & Route & 구조 \\
\midrule
\endhead
100~MHz memory address/control & 9.077~ns & 1.681~ns & 7.396~ns & 8 LUT, route 81.5\% \\
150~MHz subtract/correction 실험 & 6.677~ns & 2.896~ns & 3.781~ns & 13 levels, CARRY4 11 \\
200~MHz DSP product 결합 실험 & 8.168~ns & 4.600~ns & 3.568~ns & 17 levels, CARRY4 12 \\
최종 150~MHz arithmetic path & 6.522~ns & 2.494~ns & 4.028~ns & 7 levels, CARRY4 4 \\
최종 complete BRAM system path & 6.550~ns & 2.920~ns & 3.630~ns & 10 levels, CARRY4 6 \\
\bottomrule
\end{longtable}

여기서 얻는 직관은 다음과 같다.

\begin{itemize}
  \item LUT level이 8개여도 route가 7.396 ns이면 Boolean 식만 줄여서는 부족하다.
  \item CARRY4가 11--12개 이어진 연산은 150/200 MHz 한 cycle에 두기 어렵다.
  \item carry를 4--6개로 줄여도 3.6--4.0 ns route가 붙으므로 150 MHz가 빠듯하다.
  \item path의 이름이 add, MUX 또는 BRAM address라는 사실보다 그 앞뒤에 붙은
        selection과 physical distance가 중요하다.
\end{itemize}

\section{Adder와 correction의 timing}

23-bit add 하나는 isolated synthesis에서 CARRY4 여섯 구간이었다. dedicated
carry가 빠르더라도 comparator, second add/subtract와 MUX가 같은 cycle에 붙으면
carry가 두 번 직렬로 갈 수 있다.

\begin{lstlisting}
// Two arithmetic decisions may occupy one register boundary.
wire [23:0] diff = {1'b0, a} - {1'b0, b};
wire [23:0] corr = diff + q;
assign y = diff[23] ? corr[22:0] : diff[22:0];
\end{lstlisting}

range를 이용해 sign/borrow만으로 correction 조건을 만들 수 있는지, raw result와
correction을 서로 다른 stage에 둘 수 있는지 검토한다. pipeline을 나눌 때에는
result만 움직이지 말고 valid, mode와 writeback tag도 같이 옮긴다.

\section{MUX가 arithmetic 앞과 뒤에 있을 때}

\begin{lstlisting}
// MUX then ADD
assign y0 = a + (sel ? b : c);

// ADD in parallel, then MUX
assign y1 = sel ? (a + b) : (a + c);
\end{lstlisting}

첫 구조의 critical path는 select/data MUX 뒤에 carry가 시작될 수 있다. 두 번째
구조는 두 add가 병렬로 진행된 뒤 result MUX를 거치므로 timing이 나아질 수 있지만
adder가 두 개로 늘 수 있다. 합성기는 Boolean/arithmetic equivalence를 이용해
다시 구조를 바꿀 수 있으므로 RTL 모양만으로 단정하지 않는다.

Unified NTT의 공유 multiplier에서는 DSP 수 절감 가치가 크기 때문에 operand
selection을 허용하되 mode를 transaction 동안 고정하고 stage boundary에 맞췄다.
반면 butterfly의 add/subtract result를 너무 일찍 한 rail로 합치면 carry 뒤에
mode MUX가 붙어 150 MHz 실패 경로가 되었다. 같은 ``공유''라도 operator 크기와
MUX 위치에 따라 결론이 달랐다.

\section{Fanout과 route delay}

fanout은 한 net이 구동하는 load 수다. mode, reset, enable, state decode처럼 넓은
hierarchy를 가로지르는 control은 작은 Boolean logic이어도 많은 sink와 긴 route를
만든다.

\begin{lstlisting}
wire do_core = busy & mode_ok & phase_core;
// do_core drives many datapath, memory and valid registers.
\end{lstlisting}

해결 후보는 기능에 따라 다르다.

\begin{itemize}
  \item decode를 source 가까이에서 한 번 계산하고 register한다.
  \item 물리적으로 먼 region마다 동일한 control register를 복제한다.
  \item phase별 data command를 먼저 register해 downstream MUX의 control 수를 줄인다.
  \item 이미 FSM이 보장하는 중복 guard를 assertion과 함께 제거한다.
\end{itemize}

register replication은 FF를 조금 늘리고 route를 줄이는 trade-off다. 합성기가
동일 register를 합칠 수 있으므로 필요하면 keep/max-fanout constraint를 검토하고,
실제 replicated cell과 path를 report에서 확인한다.

\section{BRAM과 DSP 경계의 timing}

BRAM은 보통 synchronous read를 사용한다. address와 enable이 clock edge 전에
안정되어야 하고 data는 정해진 read latency 뒤에 나온다. BRAM pin 바로 앞에
bank decode, address transform과 phase MUX를 모두 놓으면 memory input path가
critical해질 수 있다.

\begin{lstlisting}
// First register a physical command.
always @(posedge clk) begin
  mem_addr_q <= mapped_addr;
  mem_bank_q <= mapped_bank;
  mem_we_q   <= request_we;
end
// The BRAM-facing MUX now sees registered fields.
\end{lstlisting}

DSP도 multiplier 앞뒤 register가 hard block 안으로 들어가는지에 따라 path가
달라진다. RTL에 wire 이름을 하나 더 붙이는 것은 pipeline이 아니다. synthesis
결과에서 DSP A/B/M/P register와 path start/end pin을 확인한다.

\section{Pipeline stage를 어디에 넣을 것인가}

긴 식의 정확히 가운데 bit 수만 보고 자르는 것보다 delay가 큰 구조 사이의
경계를 고른다.

\begin{itemize}
  \item MUX $\rightarrow$ carry가 길면 operand selection 뒤를 register한다.
  \item carry $\rightarrow$ correction carry가 길면 raw result를 register한다.
  \item DSP output $\rightarrow$ wide add가 길면 DSP PREG 또는 output stage를 사용한다.
  \item address mapping $\rightarrow$ BRAM pin이 길면 physical command를 register한다.
  \item high-fanout decode가 길면 local registered control로 분산한다.
\end{itemize}

pipeline은 latency를 늘리지만 initiation interval을 반드시 늘리지는 않는다. 각
stage가 매 cycle 새 입력을 받아 독립적으로 진행하면 latency가 늘어도 throughput은
유지된다. 반면 한 operator를 여러 cycle 재사용하는 FSM은 initiation interval도
늘 수 있다.

\section{Timing 개선의 원인을 확인하는 법}

\begin{checklistbox}
\begin{enumerate}
  \item 변경 전후 같은 clock constraint와 implementation 범위를 사용한다.
  \item startpoint, endpoint, data path delay와 slack을 기록한다.
  \item logic delay, route delay, logic level, CARRY4 수와 fanout을 분리한다.
  \item register를 추가했다면 모든 data/control branch의 cycle alignment를 검증한다.
  \item 한 경로를 고친 뒤 새 worst path가 어디로 이동했는지 확인한다.
  \item directive만 바꾼 결과와 RTL 구조를 바꾼 결과를 구분한다.
  \item OOC internal path와 complete-system endpoint path를 섞지 않는다.
\end{enumerate}
\end{checklistbox}
